\documentclass[11pt]{amsart}
\usepackage[margin=1.1in]{geometry}
\usepackage{amsmath,amssymb,amsthm}
\usepackage{booktabs}
\usepackage[colorlinks=true,linkcolor=blue,citecolor=blue,urlcolor=blue]{hyperref}

\newtheorem{theorem}{Theorem}[section]
\newtheorem{lemma}[theorem]{Lemma}
\newtheorem{proposition}[theorem]{Proposition}
\newtheorem{corollary}[theorem]{Corollary}
\theoremstyle{definition}
\newtheorem{definition}[theorem]{Definition}
\newtheorem{example}[theorem]{Example}
\newtheorem{axiom}{Axiom}
\theoremstyle{remark}
\newtheorem{remark}[theorem]{Remark}

\newcommand{\Hil}{\mathcal H}
\newcommand{\eps}{\varepsilon}
\newcommand{\one}{\mathbf 1}

\newcommand{\tr}{\operatorname{tr}}
\newcommand{\Tr}{\operatorname{Tr}}
\newcommand{\sinc}{\operatorname{sinc}}
\newcommand{\Desc}{\mathcal D}
\emergencystretch=3em

\begin{document}

\title[The Feynman path sum at $\alpha=2$ of the CAL]{An Axiomatic Derivation of the Feynman Path Sum at the Light-Cone Kernel Width and Closure Order $\alpha=2$ of the Causal Accumulation Law}

\author{Sam Leizerman}
\address{LIOR Research Initiative}
\email{samleizerman@outlook.com}
\thanks{ORCID: 0009-0000-0133-2291. Companion to \emph{The Causal Accumulation Law: The Four Corners of $\alpha$, Algebraic Closure, and the Standard Model}~\cite{CAL}, to \emph{Schr\"odinger evolution at the $\alpha\to1^-$ edge of the Causal Accumulation Law}~\cite{Edge}, and to \emph{An Axiomatic Derivation of the Born Rule at the $\delta$-Kernel Limit and Closure Order $\alpha=2$}~\cite{Born}. Dated November 12, 2026.}

\begin{abstract}
We take the Causal Accumulation Law (CAL) at closure order $\alpha=2$ and give its memory the light-cone width: at lag $\tau$ the spatial profile of the kernel is the light front of radius $c\tau$. Four results follow.

\emph{Lorentz invariance selects the order.} The symbol of the order-$\alpha$ light-cone memory is $(\tau_cR)^{-\alpha}\sin((\alpha-1)\varphi)/((\alpha-1)\sin\varphi)$, with $s=R\cos\varphi$ and $c\lvert k\rvert=R\sin\varphi$. It depends on $s^2+c^2k^2$ alone exactly when $\alpha=2$ (and trivially at $\alpha=0$); the leading defect is $\alpha(2-\alpha)\varphi^2/6$.
\begin{itemize}
\item At $\alpha=2$ the kernel is the retarded Green's function of the wave operator: Kirchhoff's shell in $3+1$ dimensions.
\item In $1+1$ dimensions it is the product of one order-one accumulation along each null leg. The order two is one plus one.
\end{itemize}

\emph{The $\alpha=2$ field equation is Klein--Gordon.} With coincidence weight $r_1>0$, the law $T=2J+r_1K_2J$ gives $(\square+\mu^2)J=0$ once the source is off, with $\mu^2=r_1/(2\tau_c^2)$, so $mc^2=\hbar\sqrt{r_1/2}/\tau_c$. This holds at finite parameters; no double scaling is needed. Carrying the channel phase of the $\alpha\to1^-$ edge to $\alpha=2$ would multiply the memory by $-1$ and make the mass tachyonic, so we use the additive law.

\emph{The field equation is a path sum.} Its Neumann series converges absolutely for every bounded source with zero past. It is a sum over piecewise-null histories: $n+1$ light-cone legs, with the real amplitude $-\mu^2$ at each of the $n$ corners. It resums to the retarded massive propagator.
\begin{itemize}
\item In $1+1$ dimensions the $n$-th term equals the $(2n+1)$-corner amplitude of Feynman's checkerboard divided by $2ic\mu$, because $(-\mu^2)^n=(i\mu)^{2n}$.
\item In $3+1$ dimensions, by descent, the histories without corners live on the cone, and the corners alone fill its interior. The Bessel tail of the massive propagator is the sum over histories with at least one corner.
\end{itemize}

\emph{Unitarity, the readout, and Feynman's kernel.}
\begin{itemize}
\item Each sign of frequency evolves by a unitary group generated by $\sqrt{m^2c^4-\hbar^2c^2\nabla^2}$.
\item On the positive-frequency sector the trace of the carrier closure, $\int JJ^\dagger$, is conserved, so the trace readout keeps its normalization in every frame. On mixed frequencies it does not.
\item Removing the rest phase and letting $c\to\infty$, the positive-frequency part of the light-cone path sum converges to Feynman's free kernel with error $O(1/c^2)$, and its composition law is Feynman's time slicing.
\end{itemize}
The $\alpha=2$ dynamics is real. The $i$ of Feynman's kernel enters only through the choice of one sign of frequency.

Numerical runs confirm each statement.
\end{abstract}

\maketitle
\tableofcontents

%======================================================================
\section{Scope and claim}

The companion paper~\cite{Edge} obtains the free Schr\"odinger equation at the $\alpha\to1^-$ edge of the CAL. There the causal structure is time ordering alone, the memory is a one-way accumulator, and nothing interacts. The Lorentz signature, and with it the light cone, belongs to the quadratic closure $JJ^\dagger$ at $\alpha=2$~\cite{CAL,Edge}. This paper takes the CAL at $\alpha=2$, with the memory spread over the light cone, and asks what its field equation is.

\begin{quote}\itshape
Give the order-$\alpha$ memory the light-cone width. Its symbol is Lorentz invariant exactly at $\alpha=2$ (Theorem~\ref{thm:lorentz}), where the kernel is the retarded Green's function of $\square$ (Corollary~\ref{cor:green}). The $\alpha=2$ field equation is then the Klein--Gordon equation, at finite parameters (Theorem~\ref{thm:KG}). Its Neumann series is a sum over piecewise-null histories with amplitude $-\mu^2$ per corner. It converges absolutely and resums to the retarded massive propagator (Theorem~\ref{thm:pathsum}). In $1+1$ dimensions it is Feynman's checkerboard term by term (Proposition~\ref{prop:checker}). In $3+1$ dimensions the interior of the cone is filled by the corners alone (Proposition~\ref{prop:descent3}). The past enters only through the pair $(J,\partial_tJ)$ (Proposition~\ref{prop:pair}). Each sign of frequency is unitary, and on positive frequency the carrier trace is conserved (Theorem~\ref{thm:unitary}). As $c\to\infty$ the positive-frequency path sum converges to Feynman's free kernel (Theorem~\ref{thm:NR}, Corollary~\ref{cor:feynman}).
\end{quote}

\begin{remark}[Bookkeeping]\label{rem:ledger}
\emph{Inputs:}
\begin{itemize}
\item the CAL (Axiom~\ref{ax:CAL});
\item closure order $\alpha=2$ with the light-cone profile, on the gauge-neutral sector (Axiom~\ref{ax:order2});
\item the additive field equation with a source of zero past, without the channel phase (Axiom~\ref{ax:FE}, Remark~\ref{rem:phase});
\item for the readout only, the restriction to one sign of frequency (Theorem~\ref{thm:unitary}).
\end{itemize}
\emph{Derived:} every proposition, theorem and corollary of \S\S4--8.

\emph{Not claimed:}
\begin{itemize}
\item interaction between two carriers at $\alpha=2$;
\item spin: each entry of $J$ obeys the scalar Klein--Gordon equation here, and a first-order (Dirac) form is not derived;
\item a value for the mass, or its relation to the edge mass of~\cite{Edge};
\item a frame-independent position readout.
\end{itemize}
\end{remark}

The $1+1$ checkerboard and its nonrelativistic limit are known~\cite{FH,JS,GJKS}, and so are chain sums for the retarded propagator on causal sets~\cite{Johnston}. What is new here is the source: the history sum is the memory series of the CAL at $\alpha=2$, and the light-cone width is selected by Lorentz invariance of the memory symbol.

%======================================================================
\section{Definitions}

Throughout, $d\in\{1,2,3\}$ is the spatial dimension, $s$ is a complex frequency with ${\rm Re}\,s>0$, $k\in\mathbb R^d$ is a wave vector, $\tau_c$ is the memory time unit, and $E_c=\hbar/\tau_c$. Complex powers take the principal branch.

\begin{definition}[Light-cone profile]\label{def:profile}
For $\tau>0$, the light-cone profile $\nu_\tau$ is the probability measure on $\mathbb R^d$ that is rotation invariant, supported in $\lvert y\rvert\le c\tau$, and has Fourier transform
\[
\hat\nu_\tau(k)=\int e^{-ik\cdot y}\,\nu_\tau(dy)=\sinc(c\lvert k\rvert\tau),\qquad \sinc z=\frac{\sin z}{z}.
\]
Explicitly:
\begin{itemize}
\item $d=3$: the uniform measure on the sphere $\lvert y\rvert=c\tau$, i.e.\ $\delta(\lvert y\rvert-c\tau)/(4\pi c^2\tau^2)$;
\item $d=2$: the density $\theta(c\tau-\lvert y\rvert)\big/\big(2\pi c\tau\sqrt{c^2\tau^2-\lvert y\rvert^2}\big)$;
\item $d=1$: the uniform density $\theta(c\tau-\lvert y\rvert)/(2c\tau)$.
\end{itemize}
The profiles for $d=2$ and $d=1$ are the orthogonal projections of the sphere in $d=3$. For $d=1$ this is Archimedes' hat-box theorem. Projection restricts the Fourier transform to a subspace, so all three have the same transform.

The \emph{light-cone width} is the statement that the profile at lag $\tau$ has radius $c\tau$. By contrast, the profile at the edge~\cite[Def.~2.2]{Edge} does not depend on the lag.
\end{definition}

\begin{definition}[Order-$\alpha$ light-cone memory]\label{def:memory}
For $\alpha>0$ and $f$ a function on spacetime,
\[
(K_\alpha f)(t,x)=\tau_c^{-\alpha}\int_0^\infty\frac{\tau^{\alpha-1}}{\Gamma(\alpha)}\int_{\mathbb R^d}f(t-\tau,x-y)\,\nu_\tau(dy)\,d\tau .
\]
It integrates over the closed past light cone $J^-(t,x)=\{(t',x'):\lvert x-x'\rvert\le c(t-t')\}$, and in $d=3$ over the cone itself.

A function $f$ has \emph{zero past} if $f=0$ for $t<t_0$, for some $t_0$. For bounded $f$ with zero past, $K_\alpha f$ is finite and
\[
\sup_x\lvert K_\alpha f(t,x)\rvert\le\tau_c^{-\alpha}\frac{(t-t_0)^\alpha}{\Gamma(\alpha+1)}\sup\lvert f\rvert,
\]
because $\nu_\tau$ is a probability measure.
\end{definition}

\begin{definition}[Symbol]\label{def:symbol}
On $e^{st+ik\cdot x}$ with ${\rm Re}\,s>0$, $K_\alpha$ acts by multiplication by
\[
m_\alpha(s,k)=\tau_c^{-\alpha}\int_0^\infty\frac{\tau^{\alpha-1}}{\Gamma(\alpha)}\,e^{-s\tau}\,\sinc(c\lvert k\rvert\tau)\,d\tau .
\]
On a carrier $e^{-iEt/\hbar+ik\cdot x}$ with real $E$, the value is the limit ${\rm Re}\,s\downarrow0$ at $s=-iE/\hbar$, as for the $\alpha=1$ value in~\cite[Def.~2.1]{Edge}.

For real $s>0$ and $k\ne0$ write
\[
s=R\cos\varphi,\qquad c\lvert k\rvert=R\sin\varphi,\qquad R=\sqrt{s^2+c^2k^2},\quad\varphi\in(0,\tfrac\pi2).
\]
\end{definition}

The wave operator is $\square=\partial_t^2-c^2\nabla^2$, with symbol $s^2+c^2k^2$.

%======================================================================
\section{Axioms and the field equation}

\begin{axiom}[Causal accumulation law; C8 of~\cite{CAL}]\label{ax:CAL}
\[
T(x)=\alpha J(x)-(1-\alpha)\int_{J^-(x)}k(\tau;x,x')\,\mathcal P(x,x')\,J(x')\,d^4x'.
\]
The kernel $k$ and the transport $\mathcal P$ are real, and $J$ takes values in $\mathbb C\otimes\mathbb H\cong M_2(\mathbb C)$.
\end{axiom}

\begin{axiom}[Closure order two at the light-cone width]\label{ax:order2}
$\alpha=2$. The kernel is $k(\tau,y)=\tau_c^{-2}\,\tau\,\nu_\tau(y)$ (Definitions~\ref{def:profile} and~\ref{def:memory}). On the gauge-neutral sector the transport is trivial, and the projected memory is $r_1K_2$, where the real number $r_1$ is the coincidence weight of the fiber factors, as in~\cite[Def.~2.4]{Edge}.
\end{axiom}

At $\alpha=2$ the prefactor $-(1-\alpha)$ is $+1$, so the law reads $T=2J+r_1K_2J$.

\begin{axiom}[Field equation]\label{ax:FE}
For a bounded source $S$ with zero past,
\begin{equation}\label{eq:FE}
T=2J+r_1K_2J=S .
\end{equation}
Wherever $S=0$ the equation is $T=0$. The solutions considered have zero past.
\end{axiom}

A field is started by a source acting on a zero past, as in~\cite{Edge}; no initial-value data are used.

\begin{remark}[The channel phase at $\alpha=2$]\label{rem:phase}
The edge paper carries the orientation $e^{-i\sigma\alpha\pi/2}$ on the memory channel~\cite[Def.~2.3]{Edge}. At $\alpha=2$ this is $e^{-i\sigma\pi}=-1$ for both orientations. Carrying it would replace $r_1$ by $-r_1$, and so $\mu^2$ by $-\mu^2$ in Theorem~\ref{thm:KG}. The dispersion would be $E^2=\hbar^2c^2k^2-m^2c^4$: modes with $c\lvert k\rvert<\mu$ grow at rate $\sqrt{\mu^2-c^2k^2}$, and there is no bounded causal response. This is the same obstruction as~\cite[Thm.~5.8]{Edge}. Table~\ref{tab:disp} shows both signs.

Only the product of the channel sign and $r_1$ enters. A real mass needs that product to be positive. Axiom~\ref{ax:FE} takes the additive law, without the phase and with $r_1>0$. If the phase is kept at $\alpha=2$, the coincidence weight must be negative there.
\end{remark}

%======================================================================
\section{The light-cone memory: symbol and Lorentz invariance}

\begin{proposition}[Symbol]\label{prop:symbol}
For $\alpha>0$, ${\rm Re}\,s>0$ and $k\ne0$, with $b=c\lvert k\rvert$,
\[
\tau_c^{\alpha}\,m_\alpha(s,k)=\frac{(s-ib)^{1-\alpha}-(s+ib)^{1-\alpha}}{2ib\,(\alpha-1)}\quad(\alpha\ne1),\qquad
\tau_c\,m_1(s,k)=\frac{\arctan(b/s)}{b}.
\]
For real $s>0$, $\tau_c^\alpha m_\alpha=R^{-\alpha}h_\alpha(\varphi)$ with
\[
h_\alpha(\varphi)=\frac{\sin((\alpha-1)\varphi)}{(\alpha-1)\sin\varphi}\quad(\alpha\ne1),\qquad h_1(\varphi)=\frac{\varphi}{\sin\varphi}.
\]
At $k=0$ the symbol is $(\tau_cs)^{-\alpha}$, the Weyl symbol of~\cite[Def.~2.1]{Edge}. The symbol does not depend on $d$.
\end{proposition}

\begin{proof}
Write $\sinc(b\tau)=\sin(b\tau)/(b\tau)$ and use $\int_0^\infty\tau^{\beta-1}e^{-z\tau}d\tau=\Gamma(\beta)z^{-\beta}$ for ${\rm Re}\,z>0$, with $\beta=\alpha-1$ and $z=s\mp ib$. The integral converges at $\tau=0$ for $\alpha>0$, because $\tau^{\alpha-2}\sin(b\tau)=O(\tau^{\alpha-1})$. Then $\Gamma(\alpha-1)/\Gamma(\alpha)=1/(\alpha-1)$ gives the first formula, and its limit $\alpha\to1$ gives the second. For real $s$, $s\mp ib=Re^{\mp i\varphi}$. The transform $\hat\nu_\tau$ is the same for every $d$ (Definition~\ref{def:profile}).
\end{proof}

\begin{theorem}[Lorentz invariance selects $\alpha=2$]\label{thm:lorentz}
For $\alpha>0$ the following are equivalent:
\begin{enumerate}
\item $m_\alpha(s,k)$ depends on $(s,k)$ only through $s^2+c^2k^2$;
\item $\alpha=2$.
\end{enumerate}
Then $m_2(s,k)=\tau_c^{-2}/(s^2+c^2k^2)$. In the limit $\alpha\to0^+$ the memory becomes the identity, which is trivially invariant. For small $\varphi$,
\[
h_\alpha(\varphi)=1+\frac{\alpha(2-\alpha)}{6}\,\varphi^2+O(\varphi^4),
\]
so the Lorentz defect is positive for $0<\alpha<2$ and negative for $\alpha>2$. The function $h_\alpha$ is symmetric under $\alpha\mapsto2-\alpha$, about the edge $\alpha=1$.
\end{theorem}

\begin{proof}
For real $s>0$ and $k\ne0$, $(R,\varphi)$ is a bijection onto $(0,\infty)\times(0,\pi/2)$. Statement~(1) holds on this set exactly when $h_\alpha$ is constant. Both sides of~(1) are analytic in $s$ for ${\rm Re}\,s>0$, so the real half-line determines the whole half-plane.

Expanding, $\sin((\alpha-1)\varphi)=(\alpha-1)\varphi\,[1-(\alpha-1)^2\varphi^2/6+O(\varphi^4)]$ and $\sin\varphi=\varphi\,[1-\varphi^2/6+O(\varphi^4)]$. Their quotient gives the expansion, which also holds at $\alpha=1$. A constant $h_\alpha$ needs $\alpha(2-\alpha)=0$, so $\alpha=2$ for $\alpha>0$. Conversely $h_2\equiv1$. Symmetry: $h_\alpha$ depends on $\alpha-1$ only through $\sin(\beta\varphi)/\beta$, which is even in $\beta$.
\end{proof}

On a carrier, $s=-iE/\hbar+0$, and $s^2+c^2k^2=-(E^2-\hbar^2c^2k^2)/\hbar^2$ with the retarded prescription. This is the Minkowski square of the four-momentum: the $\alpha=2$ memory is a function of $p^2$ and of the sign of $E$, both invariant under orthochronous Lorentz transformations.

\begin{corollary}[The order-2 kernel is the retarded Green's function of $\square$]\label{cor:green}
$\tau_c^2K_2f=G*f$ with $G(t,\cdot)=\theta(t)\,t\,\nu_t$, and $\square(G*f)=f$ for $f$ with zero past. Explicitly,
\[
G^{(3)}=\frac{\theta(t)\,\delta(\lvert x\rvert-ct)}{4\pi c^2t},\qquad
G^{(2)}=\frac{\theta(ct-\lvert x\rvert)}{2\pi c\sqrt{c^2t^2-\lvert x\rvert^2}},\qquad
G^{(1)}=\frac{\theta(ct-\lvert x\rvert)}{2c}.
\]
$G^{(3)}$ is Kirchhoff's shell~\cite{CH}.
\end{corollary}

\begin{proof}
By Theorem~\ref{thm:lorentz} the symbol of $G$ is $1/(s^2+c^2k^2)$, the inverse of the symbol of $\square$. Both $G$ and $f$ vanish in the past, so $G*f$ is the retarded solution. The explicit forms are $t\,\nu_t$ from Definition~\ref{def:profile}.
\end{proof}

\begin{proposition}[$\alpha=2$ is $1+1$: two null accumulations]\label{prop:nulllegs}
In $d=1$ let $u=ct-x$ and $v=ct+x$. Then $\square=4c^2\partial_u\partial_v$ and
\[
(G*f)(u,v)=\frac1{4c^2}\int_{-\infty}^{u}\!\!\int_{-\infty}^{v}f(u',v')\,dv'\,du' .
\]
So $G=(4c^2)^{-1}W^1_uW^1_v$, the product of one order-one causal accumulation along each null direction. The order-two memory is two order-one accumulators of the kind found at the edge~\cite[\S6]{Edge}, one per null leg.
\end{proposition}

\begin{proof}
$\partial_t=c(\partial_u+\partial_v)$ and $\partial_x=\partial_v-\partial_u$, so $\partial_t^2-c^2\partial_x^2=4c^2\partial_u\partial_v$. The Jacobian is $du\,dv=2c\,dt\,dx$, and $G^{(1)}=\theta(u)\theta(v)/(2c)$.
\end{proof}

\begin{proposition}[Descent]\label{prop:descent}
For an even function $f_1$ on $\mathbb R$, define the radial function $(\Desc f_1)(y)=-(2\pi\lvert y\rvert)^{-1}f_1'(\lvert y\rvert)$ on $\mathbb R^3$. Then $\widehat{\Desc f_1}(k)=\hat f_1(\lvert k\rvert)$. In particular $G^{(3)*n}(t,\cdot)=\Desc\,G^{(1)*n}(t,\cdot)$ for every $n\ge1$, and the same holds for any convergent series in powers of $G$.
\end{proposition}

\begin{proof}
For a radial function, $\hat f_3(\kappa)=(4\pi/\kappa)\int_0^\infty rf_3(r)\sin(\kappa r)\,dr$. With $f_3=\Desc f_1$, an integration by parts gives $2\int_0^\infty f_1(r)\cos(\kappa r)\,dr=\hat f_1(\kappa)$; for distributions it holds weakly. The spatial transform of $G^{*n}(t,\cdot)$ is the $n$-fold time convolution of $\theta(t)\sin(c\lvert k\rvert t)/(c\lvert k\rvert)$. That is the same function of $\lvert k\rvert$ for $d=1$ and $d=3$. This is the classical method of descent~\cite{CH}, read on symbols.
\end{proof}

%======================================================================
\section{The \texorpdfstring{$\alpha=2$}{alpha=2} field equation is Klein--Gordon}

\begin{theorem}[Klein--Gordon]\label{thm:KG}
Let $r_1>0$, and set
\[
\mu^2=\frac{r_1}{2\tau_c^2},\qquad mc^2=\hbar\mu=\frac{\hbar}{\tau_c}\sqrt{\frac{r_1}2}.
\]
\begin{enumerate}
\item For every bounded source $S$ with zero past, \eqref{eq:FE} has exactly one locally bounded solution with zero past. It is
\[
J=\tfrac12S-\tfrac12\mu^2\,G_\mu*S,
\]
where $G_\mu$ is the retarded fundamental solution of $\square+\mu^2$.
\item $(\square+\mu^2)J=\tfrac12\square S$ in the sense of distributions. On every open set where $S=0$, $J$ solves the Klein--Gordon equation $(\square+\mu^2)J=0$.
\item A plane wave $e^{-iEt/\hbar+ik\cdot x}$ solves $T=0$, with the memory valued as in Definition~\ref{def:symbol}, if and only if $E^2=\hbar^2c^2k^2+m^2c^4$.
\end{enumerate}
The mass is fixed at finite $r_1$ and $\tau_c$. No double scaling is needed, unlike at the edge~\cite[Def.~2.5]{Edge}.
\end{theorem}

\begin{proof}
By Corollary~\ref{cor:green}, $r_1K_2=2\mu^2G*$, so \eqref{eq:FE} is $J=\tfrac12S-\mu^2G*J$. Existence and uniqueness are Theorem~\ref{thm:pathsum}(1)--(2), and the resummation is Theorem~\ref{thm:pathsum}(3).

(2) $\square G_\mu*S=S-\mu^2G_\mu*S$, so $(\square+\mu^2)J=\tfrac12\square S+\tfrac12\mu^2S-\tfrac12\mu^2S$.

(3) The symbol of $T$ is $2+r_1\tau_c^{-2}/(s^2+c^2k^2)$. It vanishes if and only if $s^2+c^2k^2=-\mu^2$. At $s=-iE/\hbar$ this is $E^2/\hbar^2=c^2k^2+\mu^2$. On shell $E^2\ne\hbar^2c^2k^2$, so the memory value is finite.
\end{proof}

%======================================================================
\section{The path sum}

\begin{theorem}[The Neumann series is a sum over light-cone histories]\label{thm:pathsum}
Let $r_1>0$ and let $S$ be bounded with $S=0$ for $t<t_0$.
\begin{enumerate}
\item The series
\[
J=\tfrac12\sum_{n\ge0}(-\mu^2)^n\,G^{*n}*S\qquad(G^{*0}*S=S)
\]
converges absolutely and uniformly on every slab $t\le t_0+t_1$, with
\[
\sup_x\big\lvert G^{*n}*S(t,x)\big\rvert\le\frac{(t-t_0)^{2n}}{(2n)!}\,\sup\lvert S\rvert,\qquad \sup_x\lvert J(t,x)\rvert\le\tfrac12\cosh\big(\mu(t-t_0)\big)\sup\lvert S\rvert .
\]
\item Its sum is the only locally bounded solution of~\eqref{eq:FE} with zero past.
\item $G_\mu\equiv\sum_{n\ge0}(-\mu^2)^nG^{*(n+1)}$ is the retarded fundamental solution of $\square+\mu^2$, with symbol $1/(s^2+c^2k^2+\mu^2)$, and $J=\tfrac12S-\tfrac12\mu^2G_\mu*S$.
\item In $d=3$,
\[
G^{*(n+1)}(t,x)=\int\delta\Big(t-\sum_{i=0}^n\tau_i\Big)\,\delta^3\Big(x-c\sum_{i=0}^n\tau_i\hat\omega_i\Big)\prod_{i=0}^n\tau_i\,d\tau_i\,\frac{d\Omega_i}{4\pi}.
\]
Each term is an integral over histories made of $n+1$ null legs, leg $i$ of duration $\tau_i$ in direction $\hat\omega_i$, with weight $\tau_i$ per leg. So $G_\mu$ is the sum over piecewise-null histories with amplitude $-\mu^2$ at each corner.
\end{enumerate}
\end{theorem}

\begin{proof}
(1) Let $F(t)=\sup_x\lvert f(t,\cdot)\rvert$. Because $\nu_\tau$ is a probability measure, $\sup_x\lvert G*f(t,\cdot)\rvert\le\int_0^{t-t_0}\tau F(t-\tau)\,d\tau$, which is the Riemann--Liouville integral of order two of $F$ from $t_0$. By induction, the $n$-fold iterate is bounded by the order-$2n$ integral of the constant $\sup\lvert S\rvert$, which is $(t-t_0)^{2n}/(2n)!\,\sup\lvert S\rvert$. Summing $\mu^{2n}(t-t_0)^{2n}/(2n)!$ gives the cosh.

(2) Each partial sum $J_N$ satisfies $J_N=\tfrac12S-\mu^2G*J_{N-1}$; let $N\to\infty$. If two locally bounded solutions have zero past, their difference $\Delta$ satisfies $\Delta=-\mu^2G*\Delta$. Iterating $n$ times gives $\sup_x\lvert\Delta(t)\rvert\le\mu^{2n}(t-t_0)^{2n}/(2n)!\,\sup_{[t_0,t]}\lvert\Delta\rvert\to0$.

(3) $\square G^{*(n+1)}=G^{*n}$ with $G^{*0}=\delta$, so $(\square+\mu^2)G_\mu$ telescopes to $\delta$. The support is in the forward cone, and the symbol is the geometric series $\sum(-\mu^2)^n(s^2+c^2k^2)^{-n-1}$. Absolute convergence of the kernel series is shown in Propositions~\ref{prop:checker} and~\ref{prop:descent3}. Then $\tfrac12\sum_{n\ge1}(-\mu^2)^nG^{*n}*S=-\tfrac12\mu^2G_\mu*S$.

(4) $G(\tau,\cdot)=\tau\,\nu_\tau$, and $\nu_\tau$ is the uniform measure $d\Omega/4\pi$ on the sphere of radius $c\tau$. Convolving $n+1$ times gives the formula.
\end{proof}

\begin{proposition}[$1+1$: explicit terms and Feynman's checkerboard]\label{prop:checker}
Let $d=1$, $u=ct-x$, $v=ct+x$.
\begin{enumerate}
\item
\[
G^{*(n+1)}=\frac{\theta(u)\theta(v)}{2c}\,\frac1{(n!)^2}\Big(\frac{uv}{4c^2}\Big)^{n},\qquad
G_\mu=\frac{\theta(u)\theta(v)}{2c}\,J_0\Big(\mu\sqrt{t^2-x^2/c^2}\Big),
\]
and $\sum_n\mu^{2n}\lvert G^{*(n+1)}\rvert=\frac{\theta(u)\theta(v)}{2c}I_0\big(\mu\sqrt{t^2-x^2/c^2}\big)$.
\item Let $\Phi_{2n+1}(t,x)$ be the continuum checkerboard amplitude of null zigzags from $(0,0)$ to $(t,x)$ that start along one null direction and end along the other, with $2n+1$ corners and amplitude $i\mu$ per corner per unit time. Then
\[
(-\mu^2)^n\,G^{*(n+1)}=\frac{\Phi_{2n+1}}{2ic\mu}\qquad\text{for every }n\ge0,
\]
and $G_\mu=(2ic\mu)^{-1}\sum_n\Phi_{2n+1}$.
\item On a lattice of time step $\eps$, with Feynman's amplitude $i\eps\mu$ per corner~\cite{FH}, the corner-odd amplitude $L_\eps(t,x)$ satisfies $L_\eps/\eps\to\sum_n\Phi_{2n+1}=i\mu J_0\big(\mu\sqrt{t^2-x^2/c^2}\big)$, uniformly on compact subsets of the open cone.
\end{enumerate}
\end{proposition}

\begin{proof}
(1) By Proposition~\ref{prop:nulllegs}, $G*$ is $(4c^2)^{-1}$ times one integration in $u$ and one in $v$. Starting from $\theta(u)\theta(v)/(2c)$, $n$ steps give $(4c^2)^{-n}(u^n/n!)(v^n/n!)/(2c)$. The series is that of $J_0(z)=\sum_n(-1)^n(z/2)^{2n}/(n!)^2$ with $(z/2)^2=\mu^2uv/(4c^2)$, and $uv=c^2t^2-x^2$. With all signs positive it is $I_0$.

(2) Along a right-moving leg $u$ is constant and $v$ grows at rate $2c$; along a left-moving leg the roles swap. A zigzag with $2n+1$ corners has $n+1$ legs of each kind. Their durations add to $v/2c$ and $u/2c$, so the set of such zigzags has volume $(v/2c)^n/n!\cdot(u/2c)^n/n!$. Hence $\Phi_{2n+1}=(i\mu)^{2n+1}(uv/4c^2)^n/(n!)^2$, and $(i\mu)^{2n+1}/(2ic\mu)=(-\mu^2)^n/(2c)$.

(3) With $c=1$, $N=t/\eps$ steps, $N_R=(t+x)/2\eps$ right steps and $N_L=(t-x)/2\eps$ left steps,
\[
L_\eps=\sum_n(i\eps\mu)^{2n+1}\binom{N_R-1}{n}\binom{N_L-1}{n}.
\]
Here $\eps^n\binom{N_R-1}{n}\to(v/2)^n/n!$ and $\eps^n\binom{N_R-1}{n}\le(v/2)^n/n!$, so dominated convergence gives the limit. The error is $O(\eps)$ (Table~\ref{tab:checker}).
\end{proof}

\begin{remark}[Reading the checkerboard]\label{rem:checker}
The CAL insertion is real: $-\mu^2$ per corner of a light-cone history. Feynman's corner amplitude is imaginary: $i\mu$ per corner of a null zigzag. They agree term by term because $(-\mu^2)^n=(i\mu)^{2n}$. A history with $n$ insertions consists of $n+1$ steps of $G$. By Proposition~\ref{prop:nulllegs} each step is one $u$-leg and one $v$-leg, so the history has $2n+2$ legs and $2n+1$ corners: exactly the zigzags counted by $\Phi_{2n+1}$. The fixed factor $1/(2ic\mu)$ converts the corner-odd entry of the first-order checkerboard~\cite{FH,JS} into the second-order propagator $G_\mu$. In the chiral basis that entry is the mass term $\mu\,G_\mu$ of the $1+1$ Dirac propagator.
\end{remark}

\begin{proposition}[$3+1$: the corners fill the cone]\label{prop:descent3}
Let $d=3$, $r=\lvert x\rvert$ and $w=c^2t^2-r^2$. Then
\[
G^{*1}=\frac{\delta(r-ct)}{4\pi cr},\qquad G^{*(n+1)}=\theta(ct-r)\,\frac{n\,w^{n-1}}{2\pi c\,(4c^2)^n(n!)^2}\quad(n\ge1),
\]
and
\[
G_\mu=\frac{\delta(r-ct)}{4\pi cr}-\theta(ct-r)\,\frac{\mu}{4\pi c^2}\,\frac{J_1(\mu\sqrt w/c)}{\sqrt w}.
\]
The histories without corners lie on the cone (Huygens). Every history with $n\ge1$ corners lands strictly inside. So the Bessel tail of the massive propagator is the sum over histories with at least one corner. The tail series converges absolutely, with $\sum_{n\ge1}\mu^{2n}\lvert G^{*(n+1)}\rvert=\frac{\mu}{4\pi c^2}\,I_1(\mu\sqrt w/c)/\sqrt w$.
\end{proposition}

\begin{proof}
Apply $\Desc$ (Proposition~\ref{prop:descent}) to Proposition~\ref{prop:checker}(1) with $x\to r$. For $n=0$, $-\partial_r\theta(ct-r)=\delta(r-ct)$. For $n\ge1$, $\partial_rw^n=-2rnw^{n-1}$, and $w^n$ vanishes on the cone, so no shell term appears. For the sum, $J_0'=-J_1$ and $\partial_r\sqrt w=-r/\sqrt w$. The absolute bound uses $\sum_{n\ge1}nx^{n-1}/(n!)^2=I_1(2\sqrt x)/\sqrt x$. A history with legs $\tau_i\hat\omega_i$ reaches $\lvert x\rvert=c\,\lvert\sum\tau_i\hat\omega_i\rvert<c\sum\tau_i$ unless all directions coincide, a set of measure zero for $n\ge1$.
\end{proof}

%======================================================================
\section{Composition, unitarity, and the readout}

\begin{proposition}[The past enters through the pair]\label{prop:pair}
Let $S=0$ for $t\ge t_1$. For $t\ge t_1$, each spatial Fourier mode of the solution of Theorem~\ref{thm:KG} evolves by
\[
\begin{pmatrix}\hat J\\ \partial_t\hat J\end{pmatrix}(t)=M_k(t-t_1)\begin{pmatrix}\hat J\\ \partial_t\hat J\end{pmatrix}(t_1),\qquad
M_k(t)=\begin{pmatrix}\cos\omega t&\omega^{-1}\sin\omega t\\ -\omega\sin\omega t&\cos\omega t\end{pmatrix},
\]
with $\omega(k)=\sqrt{c^2k^2+\mu^2}$, and $M_k(t+t')=M_k(t)M_k(t')$.
\end{proposition}

\begin{proof}
For $t>t_1$, $J$ solves $(\square+\mu^2)J=0$ (Theorem~\ref{thm:KG}(2)), whose solutions are fixed by $(J,\partial_tJ)$ on one slice. $M_k$ is the flow of $\ddot{\hat J}=-\omega^2\hat J$.
\end{proof}

\begin{remark}[Not the faithful limit]\label{rem:faithful}
The $\alpha=2$ memory reaches over the whole past light cone, yet the past enters only through two accumulated quantities. At the edge the state was $J$ alone, kept by the fundamental theorem of calculus~\cite[Prop.~6.5]{Edge}. At $\alpha=2$ it is the pair $(J,\partial_tJ)$: in $1+1$ dimensions, one accumulator per null leg (Proposition~\ref{prop:nulllegs}). No records are kept, so this is not the faithful Markov limit, and the statement of~\cite{Edge} that a closed system cannot reach that limit is untouched.
\end{remark}

Write $\Hil=L^2(\mathbb R^d)\otimes M_2(\mathbb C)$, with the Hilbert--Schmidt product on $M_2(\mathbb C)$. Let $\hat H=\sqrt{m^2c^4-\hbar^2c^2\nabla^2}$, the self-adjoint operator of multiplication by $\hbar\omega(k)$ in Fourier space. A source-free solution has \emph{positive frequency} if every Fourier mode is proportional to $e^{-i\omega(k)t}$, and \emph{negative frequency} if every mode is proportional to $e^{+i\omega(k)t}$.

\begin{theorem}[Unitarity on each sign of frequency; the trace readout]\label{thm:unitary}
\begin{enumerate}
\item Positive-frequency solutions evolve by $J(t)=e^{-i\hat Ht/\hbar}J(0)$, a strongly continuous unitary group on $\Hil$~\cite{Stone,RS1}. Negative-frequency solutions evolve by $e^{+i\hat Ht/\hbar}$.
\item On positive frequency the closure $Q(t)=\int JJ^\dagger\,d^dx\in M_2(\mathbb C)$ is constant. So are $\Tr Q=\lVert J\rVert^2$, the readout $p_m=\Tr(P_mQ)/\Tr Q$ for projectors $P_m$ on the left factor, and the normalization of the position readout $\int_\Delta\tr(JJ^\dagger)\,d^dx/\Tr Q$.
\item On a solution with both signs of frequency, $\lVert J(t)\rVert^2$ is not constant: the cross terms oscillate at frequency $2\omega$.
\item The Klein--Gordon charge
\[
q(J)=\frac{i\hbar}{2mc^2}\int\tr\big(J^\dagger\partial_tJ-\partial_tJ^\dagger J\big)\,d^dx
\]
is conserved on every solution. It is positive on positive frequency, where $q(J)=\int(\hbar\omega/mc^2)\,\tr(\hat J^\dagger\hat J)\,d^dk$, and negative on negative frequency. On positive frequency,
\[
0\le q(J)-\lVert J\rVert^2\le\frac{\hbar^2}{2m^2c^2}\,\lVert\nabla J\rVert^2 .
\]
\end{enumerate}
\end{theorem}

\begin{proof}
(1) Each mode has the phase $e^{\mp i\omega t}$, and $\hbar\omega(k)$ is real, so this is Stone's theorem for the multiplication operator $\hat H$.

(2) $Q_{ij}=\sum_l\langle J_{jl},J_{il}\rangle_{L^2}$, and inner products are preserved by a unitary group that acts on the spatial factor only. The left projectors commute with that action.

(3) With $\hat J=a_+e^{-i\omega t}+a_-e^{i\omega t}$, $\lvert\hat J\rvert^2$ contains $2{\rm Re}\big(\bar a_+a_-e^{2i\omega t}\big)$.

(4) The charge is that of the conserved current of $(\square+\mu^2)J=0$. In Fourier space the cross terms of the two signs cancel in $q$, which leaves $\int(\hbar\omega/mc^2)(\lvert a_+\rvert^2-\lvert a_-\rvert^2)$. Finally, $1\le\hbar\omega/mc^2=\sqrt{1+\hbar^2k^2/m^2c^2}\le1+\hbar^2k^2/2m^2c^2$.
\end{proof}

\begin{remark}[Two conserved norms]\label{rem:twonorms}
On positive frequency, two positive quantities are conserved: the carrier trace $\lVert J\rVert^2=\Tr\int JJ^\dagger$, and the Klein--Gordon charge $q$.
\begin{itemize}
\item $q$ is the charge of a local conserved current, and its value is the same in every inertial frame.
\item $\lVert J\rVert^2$ is conserved in each frame, but it is not the charge of a local current, and its value depends on the frame.
\end{itemize}
They differ at order $v^2/c^2$. So the trace readout of the closure $JJ^\dagger$ is a conserved, normalized Born readout in every fixed frame, on one sign of frequency. This is consistent with~\cite[\S6.3]{Edge}: the trace is a frame-dependent density, and the Lorentz-invariant structure sits in the determinant. A frame-independent position readout is the Newton--Wigner problem~\cite{NW} and is not claimed.
\end{remark}

%======================================================================
\section{The nonrelativistic limit: Feynman's free kernel}

Fix $m$ and $\hbar$, and let $c\to\infty$; then $\mu=mc^2/\hbar\to\infty$ and $r_1=2\tau_c^2\mu^2$. Write a positive-frequency solution as $J(t)=e^{-imc^2t/\hbar}\psi_c(t)$, and let
\[
K_F(t)=e^{i\hbar t\nabla^2/2m},\qquad K_F(t;x,x')=\Big(\frac{m}{2\pi i\hbar t}\Big)^{d/2}\exp\Big(\frac{im\lvert x-x'\rvert^2}{2\hbar t}\Big)\quad(t>0),
\]
be Feynman's free kernel~\cite{Feynman,FH}.

\begin{theorem}[Nonrelativistic limit]\label{thm:NR}
For $\psi\in H^4(\mathbb R^d)\otimes M_2(\mathbb C)$,
\[
\big\lVert e^{-i(\hat H-mc^2)t/\hbar}\psi-K_F(t)\psi\big\rVert\le\frac{\lvert t\rvert\hbar^3}{8m^3c^2}\,\lVert\Delta^2\psi\rVert .
\]
Both evolutions are unitary for every $c$.
\end{theorem}

\begin{proof}
With $x=\hbar^2k^2/m^2c^2\ge0$, $-x^2/8\le\sqrt{1+x}-1-x/2\le0$. Hence $\lvert\hbar\omega-mc^2-\hbar^2k^2/2m\rvert\le\hbar^4k^4/(8m^3c^2)$. Then use $\lvert e^{ia}-e^{ib}\rvert\le\lvert a-b\rvert$ mode by mode, and Plancherel.
\end{proof}

\begin{corollary}[Feynman's kernel from the light-cone path sum]\label{cor:feynman}
Let $[G_\mu]^+$ be the positive-frequency part of $G_\mu$. In Fourier space $[\hat G_\mu]^+(t,k)=\theta(t)\,\frac{i}{2\omega}e^{-i\omega t}$, and
\[
\frac{2mc^2}{i\hbar}\,e^{imc^2t/\hbar}\,[G_\mu]^+(t)\;\longrightarrow\;K_F(t)\qquad(c\to\infty),
\]
strongly on $L^2$. For $\psi\in H^4$ the error is at most $\frac{\hbar^2}{2m^2c^2}\lVert\Delta\psi\rVert+\frac{\lvert t\rvert\hbar^3}{8m^3c^2}\lVert\Delta^2\psi\rVert$. With Theorem~\ref{thm:pathsum}, Feynman's free kernel is the limit of the sum over piecewise-null histories with amplitude $-\mu^2$ per corner,
\[
K_F(t)=\lim_{c\to\infty}\frac{2mc^2}{i\hbar}\,e^{imc^2t/\hbar}\Big[\sum_{n\ge0}(-\mu^2)^n\,G^{*(n+1)}(t)\Big]^+ .
\]
The composition $K_F(t+t')=K_F(t)K_F(t')$, which is Feynman's time slicing, is the limit of the group law of Theorem~\ref{thm:unitary}(1).
\end{corollary}

\begin{proof}
$\theta(t)\sin(\omega t)/\omega=\theta(t)\,(e^{i\omega t}-e^{-i\omega t})/2i\omega$, and its $e^{-i\omega t}$ part is $\theta(t)\,(i/2\omega)e^{-i\omega t}$. So the left side is the multiplier $(mc^2/\hbar\omega)\,e^{-i(\omega-mc^2/\hbar)t}$. Use $0\le1-mc^2/\hbar\omega\le\hbar^2k^2/2m^2c^2$ and Theorem~\ref{thm:NR}.
\end{proof}

\begin{remark}[Where the $i$ comes from]\label{rem:i}
At the edge the $i$ of the Schr\"odinger equation comes from the channel phase $e^{-i\sigma\pi/2}$~\cite{Edge}. At $\alpha=2$ the dynamics is real: $G$, the insertion $-\mu^2$ and $J_0$ are all real. The $i$ of $K_F$ enters only through the choice of one sign of frequency, which the complex carrier allows. In $1+1$ dimensions it is also the corner amplitude whose square is the real insertion (Remark~\ref{rem:checker}).
\end{remark}

\begin{remark}[Two routes to the Schr\"odinger equation]\label{rem:routes}
\begin{itemize}
\item \emph{The edge}~\cite{Edge}: $\alpha\to1^-$ at fixed memory rate, with the spatial width shrinking at fixed $\hbar^2/2m=\eps E_cr_1\langle r^2\rangle/6$. The equation is Galilean from the start, with no cone and no interaction.
\item \emph{$\alpha=2$}: at the light-cone width, on positive frequency, as $c\to\infty$ with $mc^2=\hbar\sqrt{r_1/2}/\tau_c$. As $c\to\infty$ the past light cone opens to the whole time-ordered past, which is the causal structure of the edge~\cite[\S6.3]{Edge}.
\end{itemize}
Both routes give the free Schr\"odinger equation. Whether the two masses coincide is not claimed: the coincidence weight at the two closure orders need not be the same number.
\end{remark}

%======================================================================
\section{Numerical verification}\label{sec:num}

All numbers come from \texttt{scripts/cal\_feynman\_alpha2.py} (numpy only, about 7 seconds). Units: $\hbar=c=\tau_c=1$, and $\mu=1$ unless stated.

\emph{Symbol and Lorentz invariance} (Proposition~\ref{prop:symbol}, Theorem~\ref{thm:lorentz}). Table~\ref{tab:symbol} compares direct integration with the closed form at $(s,c\lvert k\rvert)=(0.8,0.6)$, where $R=1$, and lists $h_\alpha$ at four angles. The maximal Lorentz defect $\max_\varphi\lvert h_\alpha(\varphi)-1\rvert$ over $(0,\pi/2]$ is
\[
\begin{array}{c|ccccccc}
\alpha&0&0.5&1.5&1.9&2&2.1&3\\\hline
\text{defect}&0&4.1\times10^{-1}&4.1\times10^{-1}&9.7\times10^{-2}&0&1.0\times10^{-1}&1.0
\end{array}
\]
The defect is zero only at $\alpha=0$ and $\alpha=2$, and is symmetric about $\alpha=1$.

\begin{table}[h]
\centering
\begin{tabular}{ccccccc}
\toprule
$\alpha$ & numeric & closed form & $h_\alpha(0.2)$ & $h_\alpha(0.6)$ & $h_\alpha(1.0)$ & $h_\alpha(1.4)$\\
\midrule
1.5 & 1.054093 & 1.054093 & 1.0050 & 1.0468 & 1.1395 & 1.3075\\
2.0 & 1.000000 & 1.000000 & 1.0000 & 1.0000 & 1.0000 & 1.0000\\
2.5 & 0.913547 & 0.913547 & 0.9917 & 0.9249 & 0.7903 & 0.5840\\
3.0 & 0.800000 & 0.800000 & 0.9801 & 0.8253 & 0.5403 & 0.1700\\
\bottomrule
\end{tabular}
\caption{The light-cone memory symbol: direct integration against the closed form, and the angular factor.}\label{tab:symbol}
\end{table}

\emph{Dispersion} (Theorem~\ref{thm:KG}(3), Remark~\ref{rem:phase}). With $r_1=2$, so $\mu^2=1$, the zero of the retarded symbol of $T$ is found by bisection.

\begin{table}[h]
\centering
\begin{tabular}{cccc}
\toprule
$k$ & root $E$ (additive law) & $\sqrt{k^2+\mu^2}$ & $E^2$ with the channel phase\\
\midrule
0.0 & 1.0000000000 & 1.0000000000 & $-1.00$\\
0.5 & 1.1180339887 & 1.1180339887 & $-0.75$\\
1.0 & 1.4142135624 & 1.4142135624 & $\phantom{-}0.00$\\
2.0 & 2.2360679775 & 2.2360679775 & $\phantom{-}3.00$\\
4.0 & 4.1231056256 & 4.1231056256 & $\phantom{-}15.00$\\
\bottomrule
\end{tabular}
\caption{Klein--Gordon dispersion of the $\alpha=2$ law. With the channel phase the mass is tachyonic.}\label{tab:disp}
\end{table}

\emph{Path sum in $1+1$} (Proposition~\ref{prop:checker}(1)). The terms $G^{*(n+1)}$ are built by iterated trapezoid accumulation along $u$ and then $v$, one insertion per step, on $u,v\in[0,12]$. The partial sums are compared with $\tfrac12J_0(\sqrt{uv})$.

\begin{table}[h]
\centering
\begin{tabular}{ccccccc}
\toprule
grid $h$ & $n=2$ & $n=5$ & $n=10$ & $n=20$ & $n=40$ & $n=60$\\
\midrule
0.020 & $1.44\times10^{2}$ & $1.15\times10^{3}$ & $3.28\times10^{1}$ & $3.35\times10^{-5}$ & $3.35\times10^{-5}$ & $3.35\times10^{-5}$\\
0.010 & $1.44\times10^{2}$ & $1.15\times10^{3}$ & $3.28\times10^{1}$ & $8.38\times10^{-6}$ & $8.38\times10^{-6}$ & $8.38\times10^{-6}$\\
\bottomrule
\end{tabular}
\caption{Maximum error of the partial path sums. The large intermediate values are those of the $J_0$ power series at $z=12$. The residual falls by $4.0$ when $h$ halves, so it is discretization, not the series.}\label{tab:pathsum}
\end{table}

\emph{Feynman's checkerboard} (Proposition~\ref{prop:checker}(3)). The lattice amplitude is propagated with $1$ for keeping direction and $i\eps\mu$ for a corner. The first leg moves right, and the amplitude arriving left-moving at $t=8$ is compared with $i\eps\mu\,J_0$ on $\lvert x\rvert\le7.5$.

\begin{table}[h]
\centering
\begin{tabular}{ccc}
\toprule
$\eps$ & $\max\lvert L_\eps/(i\eps\mu)-J_0\rvert$ & ratio to previous\\
\midrule
0.02 & $3.697\times10^{-2}$ & --\\
0.01 & $1.810\times10^{-2}$ & 2.04\\
0.005 & $8.950\times10^{-3}$ & 2.02\\
0.0025 & $4.451\times10^{-3}$ & 2.01\\
\bottomrule
\end{tabular}
\caption{The checkerboard's corner-odd amplitude converges to $2c\,i\mu\,G_\mu=i\mu J_0$ at first order in $\eps$.}\label{tab:checker}
\end{table}

\emph{Descent to $3+1$} (Proposition~\ref{prop:descent3}). The $1+1$ sum with $n\le60$, minus its $n=0$ term, is differentiated on the $(u,v)$ grid ($h=0.01$), and $-\partial_r/(2\pi r)$ is applied with $\partial_r=\partial_v-\partial_u$. The result is compared with $-\mu J_1(\mu\sqrt w)/(4\pi\sqrt w)$ at points with $r\ge0.25$ and $t-r\ge0.25$.

\begin{table}[h]
\centering
\begin{tabular}{cccc}
\toprule
$t$ & max error & $\max\lvert\text{exact}\rvert$ & points\\
\midrule
4 & $5.28\times10^{-7}$ & $3.090\times10^{-2}$ & 351\\
8 & $2.46\times10^{-6}$ & $2.330\times10^{-3}$ & 376\\
\bottomrule
\end{tabular}
\caption{The interior tail of the $3+1$ massive propagator, built from the summed histories.}\label{tab:descent}
\end{table}

\emph{Composition, unitarity and the nonrelativistic limit} (Proposition~\ref{prop:pair}, Theorems~\ref{thm:unitary} and~\ref{thm:NR}).
\begin{itemize}
\item Pair composition: $\max\lVert M_k(t_1+t_2)-M_k(t_2)M_k(t_1)\rVert=1.1\times10^{-14}$ over 200 random modes, with $t_1=0.7$ and $t_2=1.9$.
\item Packet $\psi\propto e^{-x^2/8+ix}$ evolved to $t=10$ by $e^{-i(\hat H-mc^2)t/\hbar}$ and by $K_F$, with $\hbar=m=1$ (Table~\ref{tab:NR}). The predicted limit of the error times $c^2$ is $(t/8)\lVert k^4\hat\psi\rVert=3.672$.
\item Trace readout: over $t\in[0,10]$ the relative spread of $\lVert J(t)\rVert^2$ is $2.5\times10^{-16}$ on positive frequency and $1.6\times10^{-1}$ on a mixed solution. The relative spread of the Klein--Gordon charge of the mixed solution is $4.4\times10^{-16}$.
\end{itemize}

\begin{table}[h]
\centering
\begin{tabular}{cccc}
\toprule
$c$ & $\lVert\text{KG}-\text{Schr\"odinger}\rVert$ & error $\times\,c^2$ & norm drift\\
\midrule
2 & $6.058\times10^{-1}$ & 2.423 & 0\\
4 & $2.076\times10^{-1}$ & 3.321 & 0\\
8 & $5.608\times10^{-2}$ & 3.589 & $2.2\times10^{-16}$\\
16 & $1.426\times10^{-2}$ & 3.651 & $2.2\times10^{-16}$\\
32 & $3.581\times10^{-3}$ & 3.667 & $2.2\times10^{-16}$\\
\bottomrule
\end{tabular}
\caption{Positive-frequency Klein--Gordon evolution, rest phase removed, against Feynman's free kernel.}\label{tab:NR}
\end{table}

%======================================================================
\section{Remarks on scope}

\begin{remark}[What is derived and what is not]\label{rem:scope}
\begin{itemize}
\item \emph{Derived from the axioms:} Lorentz invariance of the memory exactly at $\alpha=2$; the Klein--Gordon equation and the mass formula; absolute convergence of the history sum and its resummation; the checkerboard in $1+1$; the corner tail in $3+1$; the pair composition.
\item \emph{Derived under the one-sign restriction:} unitarity in $\Hil$, conservation of the closure $\int JJ^\dagger$ and of its trace readout, and the limit to Feynman's kernel.
\item \emph{Choice:} the additive law without the channel phase (Remark~\ref{rem:phase}).
\item \emph{Open:}
\begin{itemize}
\item what selects one sign of frequency;
\item interaction between carriers at $\alpha=2$, where the reciprocal (return-leg) structure that the edge lacks is available~\cite[Rem.~6.6]{Edge};
\item a first-order (spin) form in $3+1$ dimensions;
\item the relation between the $\alpha=2$ and edge masses;
\item a frame-independent position readout.
\end{itemize}
\end{itemize}
\end{remark}

\begin{remark}[Histories are memory]
Each term of the Feynman sum is the past re-entering the present $n$ times along the light cone, weighted by the order-two lag $\tau$ per leg and by $-\mu^2$ per re-entry. The sum over histories is the memory series of the CAL at $\alpha=2$. Its Markov form is the pair $(J,\partial_tJ)$, and its Galilean limit is Feynman's kernel.
\end{remark}

\appendix
%======================================================================
\section{Notation}

\begin{tabular}{ll}
$\alpha$ & closure order\\
$\nu_\tau$ & light-cone profile at lag $\tau$, radius $c\tau$ (Definition~\ref{def:profile})\\
$K_\alpha$, $m_\alpha$ & order-$\alpha$ light-cone memory and its symbol\\
$R$, $\varphi$, $h_\alpha$ & polar form of $(s,c\lvert k\rvert)$, and angular factor\\
$\square=\partial_t^2-c^2\nabla^2$ & wave operator\\
$G$, $G_\mu$ & retarded Green's functions of $\square$ and $\square+\mu^2$\\
$r_1$, $\tau_c$ & coincidence weight and memory time unit\\
$\mu^2=r_1/2\tau_c^2$, $mc^2=\hbar\mu$ & mass\\
$u=ct-x$, $v=ct+x$ & null coordinates in $1+1$\\
$\Phi_{2n+1}$ & checkerboard amplitude with $2n+1$ corners\\
$\Desc$ & descent operator (Proposition~\ref{prop:descent})\\
$w=c^2t^2-r^2$ & interval inside the cone\\
$\hat H=\sqrt{m^2c^4-\hbar^2c^2\nabla^2}$ & positive-frequency generator\\
$Q=\int JJ^\dagger$, $q$ & closure of the carrier; Klein--Gordon charge\\
$K_F$ & Feynman's free kernel\\
\end{tabular}

\section{Numerical methods}\label{app:num}

\emph{Symbol.} The integral of Definition~\ref{def:symbol} is computed on a logarithmic grid $\tau=e^u$, $u\in[-14,8]$, with $4\times10^5$ trapezoid intervals.

\emph{Dispersion.} The real part of the retarded symbol of $T$, at ${\rm Re}\,s=10^{-12}$, is bisected 200 times on $(\lvert k\rvert,50)$.

\emph{Path sums.} $G^{*(n+1)}$ is computed on the $(u,v)$ grid by cumulative trapezoid integration in $u$ and then in $v$, times $1/4$ ($c=1$). $J_0$, $J_1$ and $I$-type checks use the power series with 160 terms.

\emph{Checkerboard.} Two amplitude arrays (arriving right-moving and left-moving) are updated per time step: $R_{j+1}\leftarrow R_j+i\eps\mu L_j$ and $L_{j-1}\leftarrow L_j+i\eps\mu R_j$. Only lattice sites of the right parity are compared.

\emph{Descent.} Central differences on the $(u,v)$ grid, along the diagonal $u+v=2t$.

\emph{Packets.} A periodic grid of 4096 points on a box of length 400, with a unitary FFT, and exact spectral evolution.

%======================================================================
\begin{thebibliography}{99}
\bibitem{CAL} S.~Leizerman, \emph{The Causal Accumulation Law: The Four Corners of $\alpha$, Algebraic Closure, and the Standard Model}, LIOR Research Initiative (2026).
\bibitem{Edge} S.~Leizerman, \emph{Schr\"odinger evolution at the $\alpha\to1^-$ edge of the Causal Accumulation Law: derived unitarity, the free Galilean limit, and carrier entanglement as time dilation}, LIOR Research Initiative (2026).
\bibitem{Born} S.~Leizerman, \emph{An Axiomatic Derivation of the Born Rule at the $\delta$-Kernel Limit and Closure Order $\alpha=2$ of the Causal Accumulation Law}, LIOR Research Initiative (2026).
\bibitem{Feynman} R.~P.~Feynman, Space-time approach to non-relativistic quantum mechanics, Rev.\ Mod.\ Phys.\ \textbf{20}, 367--387 (1948).
\bibitem{FH} R.~P.~Feynman, A.~R.~Hibbs, \emph{Quantum Mechanics and Path Integrals}, McGraw-Hill (1965), Problem~2-6.
\bibitem{JS} T.~Jacobson, L.~S.~Schulman, Quantum stochastics: the passage from a relativistic to a non-relativistic path integral, J.\ Phys.\ A \textbf{17}, 375--383 (1984).
\bibitem{GJKS} B.~Gaveau, T.~Jacobson, M.~Kac, L.~S.~Schulman, Relativistic extension of the analogy between quantum mechanics and Brownian motion, Phys.\ Rev.\ Lett.\ \textbf{53}, 419--422 (1984).
\bibitem{Johnston} S.~Johnston, Particle propagators on discrete spacetime, Class.\ Quantum Grav.\ \textbf{25}, 202001 (2008).
\bibitem{CH} R.~Courant, D.~Hilbert, \emph{Methods of Mathematical Physics, Vol.~II}, Interscience (1962).
\bibitem{NW} T.~D.~Newton, E.~P.~Wigner, Localized states for elementary systems, Rev.\ Mod.\ Phys.\ \textbf{21}, 400--406 (1949).
\bibitem{Stone} M.~H.~Stone, On one-parameter unitary groups in Hilbert space, Ann.\ of Math.\ \textbf{33}, 643--648 (1932).
\bibitem{RS1} M.~Reed, B.~Simon, \emph{Methods of Modern Mathematical Physics I: Functional Analysis}, Academic Press (1980).
\end{thebibliography}

\end{document}
